\documentclass[12pt, a4paper]{article}

% 引入必要的宏包
\usepackage[UTF8]{ctex}     % 中文支持
\usepackage{amsmath, amssymb, amsthm} % 数学符号和环境
\usepackage{geometry}       % 页面设置
\usepackage{enumitem}       % 列表定制
\usepackage{fancyhdr}       % 页眉页脚
\usepackage{cases}          % 分段函数增强

% 页面边距设置
\geometry{left=2.5cm, right=2.5cm, top=2.5cm, bottom=2.5cm}

% 宏定义：方便排版选择题选项
% 横向排列的选项 (4列)
\newcommand{\choicesFour}[4]{
    \begin{enumerate}[label=\Alph*., itemsep=0pt, parsep=0pt, topsep=5pt, labelsep=0.5em]
        \begin{minipage}{0.22\linewidth} \item #1 \end{minipage}
        \begin{minipage}{0.22\linewidth} \item #2 \end{minipage}
        \begin{minipage}{0.22\linewidth} \item #3 \end{minipage}
        \begin{minipage}{0.22\linewidth} \item #4 \end{minipage}
    \end{enumerate}
}
% 横向排列的选项 (2列)
\newcommand{\choicesTwo}[4]{
    \begin{enumerate}[label=\Alph*., itemsep=0pt, parsep=0pt, topsep=5pt, labelsep=0.5em]
        \begin{minipage}{0.45\linewidth} \item #1 \end{minipage}
        \begin{minipage}{0.45\linewidth} \item #2 \end{minipage}
        \begin{minipage}{0.45\linewidth} \item #3 \end{minipage}
        \begin{minipage}{0.45\linewidth} \item #4 \end{minipage}
    \end{enumerate}
}
% 纵向排列的选项 (1列)
\newcommand{\choices}[4]{
    \begin{enumerate}[label=\Alph*., itemsep=0pt, parsep=0pt, topsep=5pt, labelsep=0.5em]
        \item #1
        \item #2
        \item #3
        \item #4
    \end{enumerate}
}

% 填空题下划线
\newcommand{\blank}[1]{\underline{\hspace{#1}}}

\title{\textbf{《高等数学》必刷习题——提高版}}
\author{}
\date{}

\begin{document}

\section*{第十一章 二重积分（提高）}

\begin{enumerate}
    \item 设函数 $ f(x,y) $ 连续，则 $ \int_{1}^{2}dy\int_{y}^{2}f(x,y)dx $ 交换积分次序为 \blank{1cm}。
    \choicesTwo
    {$ \int_{1}^{2}dx\int_{1}^{x}f(x,y)dy $}
    {$ \int_{1}^{2}dx\int_{x}^{1}f(x,y)dy $}
    {$ \int_{1}^{2}dx\int_{2}^{x}f(x,y)dy $}
    {$ \int_{1}^{2}dx\int_{x}^{2}f(x,y)dy $}

    \item 设 $D$ 是方形域： $ 0 \leq x \leq 1, 0 \leq y \leq 1, \iint_{D} xyd\sigma = $ \blank{1cm}。
    \choicesFour{1}{$ \frac{1}{2} $}{$ \frac{1}{3} $}{$ \frac{1}{4} $}

    \item $ I = \int_{0}^{1} dy \int_{-\sqrt{1-y}}^{\sqrt{1-y}} 3x^{2}y^{2} dx $ 交换积分次序为 $I =$ \blank{1cm}。
    \choicesTwo
    {$ \int_{0}^{1}dy\int_{0}^{\sqrt{1-y}}3x^{2}y^{2}dx $}
    {$ \int_{0}^{\sqrt{1-y}}dy\int_{0}^{1}3x^{2}y^{2}dx $}
    {$ \int_{-1}^{1}dx\int_{0}^{1-x^{2}}3x^{2}y^{2}dy $}
    {$ \int_{0}^{1}dx\int_{0}^{1+x^{2}}3x^{2}y^{2}dy $}

    \item 二重积分 $ \iint_{D} f(x,y)d\sigma $ 在极坐标下的面积元素为 \blank{1cm}。
    \choicesFour{$ d\sigma = dx dy $}{$ d\sigma = r dr d\theta $}{$ d\sigma = dr d\theta $}{$ d\sigma = r^{2}\sin\theta dr d\theta $}

    \item [5.] (题目序号更正) 设 $ I = \int_{0}^{1} dx \int_{x^{2}}^{x} f(x, y) dy $ 交换积分顺序后，则 $ I = $ \blank{2cm}。

    \item [6.] $D$ 是圆环 $ a^{2}\leq x^{2}+y^{2}\leq b^{2} $ ，计算 $ I=\iint_{D}(x^{2}+y^{2})d\sigma $。

    \item [7.] 交换积分顺序： $ \int_{0}^{\frac{1}{4}}dy\int_{y}^{\sqrt{y}}f(x,y)dx+\int_{\frac{1}{4}}^{\frac{1}{2}}dy\int_{y}^{\frac{1}{2}}f(x,y)dx $ 为 \blank{2cm}。

    \item [8.] 计算二重积分 $ \iint_{D} \frac{y}{x} d\sigma $ ，其中 $D$ 由 $ x^{2} + y^{2} \leq a^{2} (a > 0) $ ，$y = x$ 及 $x$ 轴在第一象限所围成的区域。

    \item [9.] 二重积分 $ I=\iint_{D}\ln(1+x^{2}+y^{2})dxdy $ ,其中 $D$ 是由圆周 $ x^{2}+y^{2}=1 $ 及坐标轴所围成的第一象限内的区域。

    \item [10.] 求二重积分 $ \iint_{D} (x + 2y) \, dx \, dy $ ，其中区域 $D$ 是由 $ y = x, y = 4x, xy = 1 $ 围成的第一象限部分的闭区域。

    \item [11.] 求 $ \iint_{D} ye^{xy} dx dy $ ，其中 $D$ 是由 $ xy = 1, x = 2, y = 1 $ 所围成的区域。

    \item [12.] 求二重积分 $ \iint_{D}\frac{y}{x}dxdy $ ，其中 $D$ 是由 $ x=e,y=\ln x,y=0 $ 所围成的区域。

    \item [13.] 计算 $ \iint_{D} \frac{x}{y} d\sigma $ ，其中 $D$ 是由 $ y = 1, y = x^{2}, x = 2 $ 所围成的闭区域。

    \item [14.] 计算二重积分 $ \iint_{D} x\sqrt{y} $ 其中 $D$ 是由 $ y = \sqrt{x}, y = x^{2} $ 围成的区域。
\end{enumerate}

\end{document}
